%═══════════════════════════════════════════════════════════════════════════════
\chapter{Demostración de la plantilla en apéndices}
\label{ap:demo}
%═══════════════════════════════════════════════════════════════════════════════

Este apéndice ilustra el uso de la plantilla \texttt{latex-catppuccin-template}
en material complementario. Todos los elementos mostrados utilizan los colores
de la paleta Catppuccin, adaptándose automáticamente al tema seleccionado.

%───────────────────────────────────────────────────────────────────────────────
\section{Entornos teoremáticos}
%───────────────────────────────────────────────────────────────────────────────

Los mismos entornos definidos para el cuerpo principal están disponibles aquí.

\begin{definicion}{Convergencia débil}
Una sucesión $\{x_n\}$ en un espacio de Hilbert $\mathcal{H}$ converge
débilmente a $x$ si $\langle x_n, y \rangle \to \langle x, y \rangle$ para todo
$y \in \mathcal{H}$.
\end{definicion}

\begin{teorema}{Banach–Alaoglu}
La bola unitaria cerrada en el dual de un espacio normado es compacta en la
topología débil-*.
\end{teorema}

\begin{proposicion}{Desigualdad de Bessel}
Sea $\{e_n\}$ una sucesión ortonormal en un espacio de Hilbert. Entonces
$\sum_n |\langle x, e_n \rangle|^2 \leq \|x\|^2$ para todo $x$.
\end{proposicion}

\begin{ejemplo}{Sucesión ortonormal en $L^2[0,1]$}
Las funciones $e_n(x) = \sqrt{2}\sin(n\pi x)$ forman una base ortonormal
para $L^2[0,1]$.
\end{ejemplo}

\begin{observacion}{Topologías débiles}
La convergencia débil es estrictamente más débil que la convergencia en norma
en espacios de dimensión infinita.
\end{observacion}

\begin{fisicaconexion}{Estados coherentes}
En óptica cuántica, los estados coherentes minimizan las relaciones de
incertidumbre y forman un sistema sobrecompleto en el espacio de Fock.
\end{fisicaconexion}

%───────────────────────────────────────────────────────────────────────────────
\section{Ecuaciones destacadas}
%───────────────────────────────────────────────────────────────────────────────

Puedes incluir ecuaciones con colores semánticos usando los comandos
\texttt{\textbackslash textcolor} con los nombres de la paleta.

\begin{equation}
\label{eq:demo2}
\textcolor{CtpBlue}{\mathcal{F}[f](\xi)} =
\textcolor{CtpPeach}{\int_{-\infty}^{\infty}}
\textcolor{CtpGreen}{f(x) e^{-2\pi i x \xi}} \, \textcolor{CtpTeal}{dx}.
\end{equation}

\begin{align}
\textcolor{CtpMauve}{\nabla \cdot \mathbf{E}} &= \textcolor{CtpSky}{\frac{\rho}{\varepsilon_0}}, \\
\textcolor{CtpMauve}{\nabla \times \mathbf{E}} &= \textcolor{CtpSky}{-\frac{\partial \mathbf{B}}{\partial t}}.
\end{align}

%───────────────────────────────────────────────────────────────────────────────
\section{Algoritmos}
%───────────────────────────────────────────────────────────────────────────────

Los algoritmos mantienen el mismo estilo que en el cuerpo principal.

\begin{algorithm}[H]
\caption{Método de Euler–Maruyama para SDEs}
\label{alg:em}
\LinesNumbered
\KwIn{Ecuación $dX_t = \mu(X_t,t)dt + \sigma(X_t,t)dW_t$,
      condición inicial $X_0$, paso $\Delta t$, tiempo final $T$}
\KwOut{Trayectoria $\{X_k\}_{k=0}^{N}$ con $N = T/\Delta t$}
$X \leftarrow X_0$\;
\For{$k = 0$ \KwTo $N-1$}{
  $t \leftarrow k\Delta t$\;
  $\xi \leftarrow \mathcal{N}(0,1)$\tcp*{ruido gaussiano}
  $X \leftarrow X + \mu(X,t)\Delta t + \sigma(X,t)\sqrt{\Delta t}\,\xi$\;
}
\Return trayectoria\;
\end{algorithm}

\begin{algorithm}[H]
\caption{Denoising score matching (un paso)}
\label{alg:dsm}
\LinesNumbered
\KwIn{Dato $\mathbf{x}_0$, tiempo $t$, red de score $\mathbf{s}_\theta$,
      schedule $\{\bar\alpha_t\}$}
\KwOut{Pérdida $\mathcal{L}$}
$\varepsilon \leftarrow \mathcal{N}(\mathbf{0}, \mathbf{I})$\;
$\mathbf{x}_t \leftarrow \sqrt{\bar\alpha_t}\,\mathbf{x}_0 + \sqrt{1-\bar\alpha_t}\,\varepsilon$\;
$\hat{\varepsilon} \leftarrow \mathbf{s}_\theta(\mathbf{x}_t, t)$\;
$\mathcal{L} \leftarrow \|\hat{\varepsilon} - \varepsilon\|^2$\;
\Return $\mathcal{L}$\;
\end{algorithm}

%───────────────────────────────────────────────────────────────────────────────
\section{Código fuente}
%───────────────────────────────────────────────────────────────────────────────

La configuración de \texttt{listings} también se aplica en los apéndices.

\begin{lstlisting}[language=Python, caption={Red neuronal con PyTorch}]
import torch
import torch.nn as nn

class ScoreNet(nn.Module):
    def __init__(self, dim, hidden=128):
        super().__init__()
        self.net = nn.Sequential(
            nn.Linear(dim, hidden),
            nn.ReLU(),
            nn.Linear(hidden, hidden),
            nn.ReLU(),
            nn.Linear(hidden, dim)
        )
    
    def forward(self, x, t):
        # t es el tiempo, se puede incorporar como embedding
        return self.net(x)
\end{lstlisting}

\begin{lstlisting}[language=julia, style=catppuccin-base, caption={Simulación de proceso OU en Julia}]
using Random, Plots

function euler_maruyama(drift, diffusion, x0, dt, T)
    N = Int(T / dt)
    x = zeros(N+1)
    x[1] = x0
    for i in 1:N
        t = (i-1)*dt
        x[i+1] = x[i] + drift(x[i], t)*dt + diffusion(x[i], t)*sqrt(dt)*randn()
    end
    return x
end

# Ejemplo: dx = -x dt + dW
drift(x, t) = -x
diffusion(x, t) = 1.0
x = euler_maruyama(drift, diffusion, 0.0, 0.01, 10.0)
plot(x)
\end{lstlisting}

%───────────────────────────────────────────────────────────────────────────────
\section{Ejercicios}
%───────────────────────────────────────────────────────────────────────────────

Los entornos de ejercicio están disponibles para plantear problemas.

\begin{ejercicio}
Demuestre que el método de Euler–Maruyama tiene orden de convergencia fuerte
$1/2$ para SDEs con coeficientes regulares.
\end{ejercicio}

\begin{ejercicio}
Verifique que la función de score para una distribución gaussiana
$\mathcal{N}(\mu, \sigma^2)$ es $\nabla_x \log p(x) = -(x-\mu)/\sigma^2$.
\end{ejercicio}

\begin{ejerciciocomp}{Implementación de DDPM}
Implemente el algoritmo de difusión denoising probabilistic model (DDPM)
para el dataset MNIST. Use una red U-Net pequeña y el schedule coseno.
Reporte muestras generadas después de 50 épocas.
\end{ejerciciocomp}

%───────────────────────────────────────────────────────────────────────────────
\section{Estrella Polar (opcional)}
%───────────────────────────────────────────────────────────────────────────────

La sección \textbf{Estrella Polar} también puede aparecer en apéndices si es
relevante.

\begin{estrella}{Espectro de potencias del reactor}
La densidad espectral de potencia de las fluctuaciones neutrónicas en un
reactor crítico puede derivarse a partir de la ecuación de Langevin. El resultado
es una lorentziana con ancho proporcional a la constante de decaimiento de los
precursores.
\end{estrella}

%───────────────────────────────────────────────────────────────────────────────
\section{Dificultad y notas}
%───────────────────────────────────────────────────────────────────────────────

Incluimos un ejemplo del comando \texttt{\textbackslash dificultad} y del
entorno \texttt{notasbib}.

\dificultad{2}{3}

\begin{notasbib}
Para una revisión exhaustiva de los métodos numéricos para SDEs, véase
\cite{kloeden1992numerical}. El algoritmo de score matching se analiza en
\cite{hyvarinen2005estimation} y su versión denoising en \cite{vincent2011connection}.
\end{notasbib}